\chapter{极限与微积分}

\section{积分定阶}
\recognize{比较无穷小的阶，式中出现短区间积分、复合上限。先把被积函数化为幂，再计算上限带来的阶数。}
\core{若$h(t)=ct^m+o(t^m)$，$m>-1$，则在相应单侧邻域有
\[\int_0^u h(t)\,dt=\frac{c}{m+1}u^{m+1}+o(u^{m+1}).\]
若$u\sim ax^r$，积分的首项是$ca^{m+1}x^{r(m+1)}/(m+1)$。$x^8$比$x^7$高阶，因为$x^8/x^7\to0$。}
\process{\begin{enumerate}\item 圈出趋近点和上限$u(x)$。\item 求被积函数首个非零项$ct^m$。\item 积分使指数加一，再代上限；保留系数。\item 两式作比，区分高阶、低阶、等价和同阶非等价。\end{enumerate}}
\fast{不必先求原函数。$e^{t^3}-1\sim t^3$直接给积分四阶；上限$x^2$再把阶数变为八阶。上限为负时保留有向积分的符号。}
\budget{易至中；选择题1—2分钟，含抵消的大题6—9分钟。}
\fullquestion{math2-2021-01}
\apply{读到$e^{t^3}-1$联想到$e^v-1\sim v$，读到$x^2$上限联想到复合定阶。于是积分$\sim x^8/4$，除以$x^7$趋于0，选C。若改问与$x^8$的关系，应答同阶但非等价；若改成$x^8/4$才是等价。}
\related{2020-01：四个短区间积分比较；2017-15：根式积分先缩放；2013-15：乘积中的二阶首项。}

\section{等价与配参}
\recognize{“与某式等价”“比$x$高阶”“求$a,b$”意味着低阶项必须消失、目标阶系数必须匹配。}
\core{\[\sin x=x-\frac{x^3}{6}+o(x^3),\quad \arcsin x=x+\frac{x^3}{6}+o(x^3),\quad
\ln(1+x)=x-\frac{x^2}{2}+\frac{x^3}{3}+o(x^3),\]
\[e^x=1+x+\frac{x^2}{2}+\frac{x^3}{6}+o(x^3),\quad(1+x)^a=1+ax+\frac{a(a-1)}2x^2+o(x^2).\]
$f\sim cx^k$要求$k$阶以下系数全为0，$k$阶系数等于$c\ne0$；仅同阶只要求该系数非零。}
\process{\begin{enumerate}\item 先确定目标阶$k$，各项展开到$k$阶。\item 按$x,x^2,\ldots$合并，先消低阶项。\item 将目标阶系数匹配；最后确认不为0。\end{enumerate}}
\fast{复合变量按最终阶数截断。例如$(1+x^2)^{1/3}-1$到$x^2$即够，$\arcsin x$也只需写$x+o(x^2)$。不必机械展开同样多项。}
\budget{易至中；选择、填空1—3分钟，大题6—8分钟。}
\fullquestion{math2-2026-01}
\apply{右端$\sqrt[3]{1+x^2}-1\sim x^2/3$先锁定二阶。左端为$(b+1)x+ax^2+o(x^2)$。先令$b+1=0$，再令$a=1/3$，选A。这是系数匹配，求导联立会徒增计算。}
\related{2023-11、2015-15：匹配到三阶；2018-01：幂指极限反推底数系数；2022-01、2025-04：用比值检验等价命题。}

\section{抵消后定阶}
\recognize{两个各自趋于无穷的量相减，或两个首项相同的函数相减。先合并，找到抵消之后真正留下的阶。}
\core{等价替换适用于乘除中的整体因子。加减式应写$f=a+\delta_1,g=a+\delta_2$，保留$\delta_1-\delta_2$。例如$\sin x\sim x$且$\ln(1+x)\sim x$，却不能推出两者之差等价于0。
\[\frac1x-\frac{\ln(1+x)}{x\sin x}=\frac{\sin x-\ln(1+x)}{x\sin x}.\]}
\process{\begin{enumerate}\item 通分、提公因子或有理化，形成一个分式。\item 由分母阶数决定分子展开精度。\item 先进行代数抵消，再取首个非零项之比。\item 若连续两阶抵消，将余项同步提高一阶。\end{enumerate}}
\fast{分母二阶，分子只需二阶；分母四阶则不能停在二阶。洛必达适合求导后明显简化的$0/0$、$\infty/\infty$；遇多层乘积和复合函数，泰勒通常更短。}
\budget{中；填空2—3分钟，大题7—10分钟。}
\fullquestion{math2-2026-12}
\apply{通分后分母$x\sin x\sim x^2$，分子$\sin x-\ln(1+x)=x^2/2+o(x^2)$，极限为$1/2$。读到两个$1/x$量级相减，就应保留下一阶；分别把$\sin x$和$\ln(1+x)$换成$x$会抹掉答案。}
\related{2021-17：积分与三角倒数相减；2012-15：极限求出常数后继续定阶；2016-15：底数前几阶抵消后再取对数。}

\section{幂指取对数}
\recognize{变量同时位于底数和指数，出现$1^\infty$、$0^0$、$\infty^0$。目标是把乘方改成指数乘对数。}
\core{底数$u(x)>0$时，$u(x)^{v(x)}=\exp\{v(x)\ln u(x)\}$。若$u\to1$且$v(u-1)\to L$，则$u^v\to e^L$，因为$\ln u\sim u-1$。若$u-1$有抵消，应先定阶。}
\process{\begin{enumerate}\item 检查底数在邻域内为正，记$L=\lim v\ln u$。\item 底数写成$1+\varepsilon$，利用$\ln(1+\varepsilon)\sim\varepsilon$。\item 求指数极限，再还原$e^L$。\end{enumerate}}
\fast{不能把“底数趋于1”直接变成答案1。若底数趋于常数$c\ne1$，先看$v\ln c$，不要套$1^\infty$公式。}
\budget{易至中；填空1—2分钟，抵消型3—5分钟。}
\fullquestion{math2-2022-11}
\apply{底数$(1+e^x)/2=1+(e^x-1)/2$，指数$\cot x\sim1/x$，所以对数极限$\lim\cot x\ln[1+(e^x-1)/2]=1/2$，原极限$\sqrt e$。触发点是“底数1加一阶小量，指数一阶倒数”。}
\related{2013-09、2019-09：底数近1；2015-02：先计算关于参数的幂指极限，再研究所得函数。}

\section{导数识别}
\recognize{极限中有$f(a+h)-f(a)$，外面乘着趋于无穷的系数，或自变量由参数方程给出。先找增量$h$，不是直接求整个表达式的导数。}
\core{\[\lim_{h\to0}\frac{f(a+h)-f(a)}h=f'(a),\qquad
\lim_{x\to\infty}x[f(a+c/x)-f(a)]=cf'(a).\]
已有可导条件可用$f(a+h)=f(a)+f'(a)h+o(h)$；“导数存在”与“邻域导函数连续”不同。}
\process{\begin{enumerate}\item 写出基点$a$与增量$h(x)$。\item 乘除$h$，分成差商与已知初等极限。\item 参数式先找到对应参数，再求$y_t/x_t$。\end{enumerate}}
\fast{极限变量可以是$n$、$1/x$或参数$t$。只要函数差与自变量差匹配，核心仍是一个导数。}
\budget{易至中；选择题1—2分钟。}
\fullquestion{math2-2024-02}
\apply{这里基点为2，增量为$2/x$，故极限为$2f'(2)$。参数式$x=1+t^3$在$x=2$对应$t=1$；$f'(2)=2e/3$，答案$4e/3$，选B。直接显式解出$f$再展开也可行，但多做了三次根式复合。}
\related{2013-02：隐函数加差商；2022-17：已知可导，复合增量均为$x^2$量级。}

\section{由极限证可导}
\recognize{题目只给$f$连续，却要求证明可导并求导数。应从条件中分离$f(x)-f(0)$的差商，不能对未知$f$先用洛必达。}
\core{证明$f'(a)=A$等价于证明$f(a+h)=f(a)+Ah+o(h)$。若已知某分式极限有限而分母趋零，先推出分子趋零；若分子中有$x f(x)$，通常还要除一次$x$才能找出$f(0)$。}
\process{\begin{enumerate}\item 只展开已知初等函数，不展开尚未证可导的$f$。\item 从给定极限解出$f(x)$的渐近式。\item 用连续性确认常数项就是$f(0)$。\item 构造$[f(x)-f(0)]/x$，取极限完成证明。\end{enumerate}}
\fast{“先证后用”：已知可导时允许写一阶展开；只知连续时，写$f=f(0)+f'(0)x+o(x)$相当于提前假定结论。}
\budget{中；10—12分题8—12分钟。}
\fullquestion{math2-2025-18}
\apply{分母$\ln(1+x)+\ln(1-x)=-x^2+o(x^2)$，给定比值为$-3$，所以分子为$3x^2+o(x^2)$。又$e^{2\sin x}=1+2x+2x^2+o(x^2)$，得到$xf(x)=2x+5x^2+o(x^2)$。于是$f(x)=2+5x+o(x)$；连续性给$f(0)=2$，差商趋5。}
\related{2022-17：已知$f$在1可导，先由有限极限得$f(1)=0$，再用两个复合差商求$f'(1)=-1$。}

\section{绝对值与可导}
\recognize{差商里出现$|f(x)|$、$|x|$，问可导的充分条件。先分$f(0)\ne0$与$f(0)=0$，前者用连续性定号，后者用左右符号约束。}
\core{若$f$在0连续且$f(0)\ne0$，则局部$|f|=\pm f$。若$f(0)=0$且$\lim_{x\to0}|f(x)|/x$存在，由右侧非负、左侧非正可知该极限只能是0，从而$|f(x)|=o(|x|)$，故$f'(0)=0$。}
\process{\begin{enumerate}\item 有限差商极限先令分子极限为0，约束$f(0)$。\item 非零值分支用连续性去绝对值。\item 零值分支比较左右符号，推出极限为0。\item 回到原函数差商，而非停在$|f|$可导。\end{enumerate}}
\fast{单独知道$|f|$可导不总能推出$f$可导；本类题中连续性负责排除零点附近的跳号反例。看到$x$而非$|x|$作分母，立即检查左右符号。}
\budget{中偏难；多命题选择题3—4分钟。}
\fullquestion{math2-2025-07}
\apply{①②均先迫使$f(0)\ge0$；若$f(0)>0$，差商直接成为$f$的差商，若为0转入③。③由左右符号推出极限0，再由夹逼得$f'(0)=0$。④在$f(0)\ne0$时局部定号，在$f(0)=0$时同③。因此四个条件都可推出可导，选D。}
\related{2018-02：根式、绝对值的局部阶数；2023-05：参数绝对值先按正负分段。}

\section{分段与振荡}
\recognize{分界点、绝对值或$x^\alpha\sin(1/x^\beta)$同时出现。分别检查连续、可导、导函数连续，不能只比较两侧导函数。}
\core{在分界点$a$依次检验：$f(a-)=f(a+)=f(a)$；左右差商相等；最后才检验$\lim f'(x)=f'(a)$。模型$f(x)=x^\alpha\sin(x^{-\beta})$（$x>0$，左侧为0，$\alpha,\beta>0$）在0可导要求$\alpha>1$，导函数连续要求$\alpha>\beta+1$。}
\process{\begin{enumerate}\item 把$|t|$或$|x|$按正负写成两支。\item 先查函数值，再用定义求分界点导数。\item 分界点外用求导规则；若问二阶，重复一次定义。\item 振荡项用“幂次趋0乘有界量”判极限，指数为0时检查振荡。\end{enumerate}}
\fast{可导不保证导函数连续。振荡函数一求导会多出$x^{-\beta-1}$，因此连续阈值通常比可导阈值更高。}
\budget{中；选择题2—4分钟。}
\fullquestion{math2-2023-05}
\apply{$t\ge0$时$x=3t,y=t\sin t$，故$f(x)=\frac x3\sin\frac x3$；$t<0$时$x=t,y=-t\sin t$，故$f(x)=-x\sin x$。两支$f(x)=O(x^2)$，所以$f'(0)=0$，且两侧$f'(x)\to0$。但$f'(x)/x$的右极限为$2/9$、左极限为$-2$，二阶导不存在，选C。}
\related{2017-01、2018-03：连续拼接；2024-01、2020-02：先找所有异常点，再逐一算左右极限。}

\section{高阶导数}
\recognize{问$f^{(n)}(0)$，或多项式乘指数的指定阶导数。原点用系数法，非零点用有限项莱布尼茨公式。}
\core{若$f(x)=\sum c_kx^k$在原点有相应泰勒展开，则$f^{(n)}(0)=n!c_n$。
\[(uv)^{(n)}=\sum_{k=0}^n\binom nk u^{(k)}v^{(n-k)}.\]
当$u$为$m$次多项式时，$k>m$项全为0；$D^n(e^{ax}P)=e^{ax}(D+a)^nP$。}
\process{\begin{enumerate}\item 原点求导先写常用展开，定位$x^n$系数。\item 若有$x^m$因子，把目标变为另一因子的$x^{n-m}$项。\item 非原点且多项式次数低，列有限个莱布尼茨项。\item 检查$n<m$、$n=m$等边界。\end{enumerate}}
\fast{问$n$阶导数不意味着连续求导$n$次。$x^2e^x$五阶只保留多项式求0、1、2阶的三项。}
\budget{易至中；选择填空1—3分钟。}
\fullquestion{math2-2024-14}
\apply{$f=x^2e^x-x^2$，五阶时$-x^2$消失。莱布尼茨得到$f^{(5)}(x)=e^x(x^2+10x+20)$，所以$f^{(5)}(1)=31e$。若问原点一般$n$阶，则改用系数法更省。}
\related{2020-04：$x^2\ln(1-x)$在0的$n$阶导数为$-n!/(n-2)$（$n\ge3$）；2015-10：$x^2 2^x$的系数平移。}

\section{参数二阶导}
\recognize{由$x=x(t),y=y(t)$求$dy/dx$或$d^2y/dx^2$。每求一次关于$x$的导数，都要除一次$x'(t)$。}
\core{当$x'(t)\ne0$，
\[\frac{dy}{dx}=\frac{y'}{x'},\qquad \frac{d^2y}{dx^2}=\frac1{x'}\frac d{dt}\!\left(\frac{y'}{x'}\right)=\frac{y''x'-y'x''}{(x')^3}.\]
参数导数上的撇号此处均表示对$t$求导。}
\process{\begin{enumerate}\item 求$x',y'$，先把$y'/x'$约简。\item 对约简结果按$t$求导，再除$x'$。\item 最后代入指定$t$；检查$x'\ne0$。\end{enumerate}}
\fast{先代参数会丢掉一阶导函数的变化信息。若$y'=q(t)x'$，立刻得到$dy/dx=q(t)$，二阶只剩$q'/x'$。}
\budget{易至中；填空2—3分钟。}
\fullquestion{math2-2021-12}
\apply{$x'=2e^t+1$，$y'=4te^t+2t=2t(2e^t+1)$。结构直接提示先约分，$y_x=2t$，$y_{xx}=2/(2e^t+1)$，在$t=0$为$2/3$。若先套完整三次分母公式，能算对但更长。}
\related{2015-09、2017-10、2020-09：同一流程；2025-14：参数中再嵌隐式变限积分。}

\section{隐函数求导}
\recognize{给$F(x,y)=0$，求一阶、二阶导数或指定点斜率。先定位$y$值，再按$y=y(x)$求导。}
\core{若$F_y\ne0$，
\[y'=-\frac{F_x}{F_y},\quad y''=-\frac{F_{xx}+2F_{xy}y'+F_{yy}(y')^2}{F_y}.\]
只求某点二阶导时，可对原式连续求导，在最后代入$y,y'$；不必解出$y(x)$。}
\process{\begin{enumerate}\item 将给定$x$代回方程，求相应$y$。\item 对原方程求导，解出该点$y'$。\item 再对一阶等式求导，注意$xy'$和$y^2y'$均用乘积法。\item 代入已知点求$y''$。\end{enumerate}}
\fast{隐函数导数的分母是$F_y$；若为0，先换成$x=x(y)$或改用参数，不强套公式。点处计算比完整化简$y''(x)$更省时。}
\budget{易至中；填空2—3分钟。}
\fullquestion{math2-2022-12}
\apply{原式$x^2+xy+y^3=3$在$x=1$给$y=1$。一阶为$2x+y+(x+3y^2)y'=0$，得$y'=-3/4$。二阶为$2+2y'+6y(y')^2+(x+3y^2)y''=0$，得$y''=-31/32$。重点是对$(x+3y^2)y'$求导保留全部乘积项。}
\related{2012-09；2023-14：隐导后取法线斜率；2026-13：隐导再代曲率公式。}

\section{反函数导数}
\recognize{函数由积分定义，却只问反函数某点导数。先找原像，再取原函数导数的倒数，无需求积分。}
\core{若$y_0=f(x_0)$且$f'(x_0)\ne0$，则
\[(f^{-1})'(y_0)=\frac1{f'(x_0)},\quad (f^{-1})''(y_0)=-\frac{f''(x_0)}{[f'(x_0)]^3}.\]
积分函数$\int_a^x h(t)dt$天然满足$f(a)=0$；若$h$在邻域内恒正或恒负，则局部单调。}
\process{\begin{enumerate}\item 解$f(x_0)=y_0$；$y_0=0$时优先看积分上下限相等。\item 用微积分基本定理求$f'$。\item 在$x_0$代值后取倒数，二阶再用上式。\end{enumerate}}
\fast{“在$y=0$处”要代入的是原像$x_0$，不是把$f'(0)$取倒数。}
\budget{易；填空1—2分钟。}
\fullquestion{math2-2013-10}
\apply{题中$f(x)=\int_{-1}^x\sqrt{1-e^t}\,dt$，$f(-1)=0$，所以原像为$-1$。$f'(-1)=\sqrt{1-e^{-1}}$，反函数导数为$1/\sqrt{1-e^{-1}}$。无需计算这个没有初等原函数的积分。}
\related{2026-06：同型题，但归档题干和网站解答存在数值不一致，详见研究附件。}

\section{变限积分求导}
\recognize{上下限含$x$，问导数、奇偶性或零点。先把积分视为原函数的复合，避免尝试计算无初等原函数的积分。}
\core{连续$h$满足
\[\frac d{dx}\int_{a(x)}^{b(x)}h(t)dt=h(b(x))b'(x)-h(a(x))a'(x).\]
若$H(x)=\int_0^x h(t)dt$，$h$偶则$H$奇；$h$奇则$H$偶。连续被积函数保证变限积分可导；仅有跳跃间断时，积分仍连续，但分界点导数需看左右差商。}
\process{\begin{enumerate}\item 分清外层$x$和积分哑变量$t$。\item 上限项减下限项，各乘其导数。\item 奇偶题先判被积函数，再判上下限复合。\item 有间断点时回到积分差商，不直接把该点函数值当导数。\end{enumerate}}
\fast{导数符号可用于单调性，积分本身往往不用算。$\int_0^x t^3dt$是偶函数，体现“奇函数积分变偶”，不要只记一个奇偶口诀。}
\budget{易至中；选择填空1—3分钟。}
\fullquestion{math2-2024-03}
\apply{$\sin(t^3)$是奇函数，令$H(u)=\int_0^u\sin(t^3)dt$，则$H$偶；$\sin x$奇，故$f(x)=H(\sin x)$偶。再积分$g(x)=\int_0^x f(t)dt$为奇，选D。无需展开任何积分。}
\related{2013-03：阶跃被积函数的积分连续不可导；2015-11：外面还有因子$x$，先乘积法；2021-14：嵌套积分逐层求导。}

\section{含参积分换元}
\recognize{积分号内出现$g(xt)$，只知$g$连续却要求求导。先换元把参数移到上限，避开不存在的$g'$。}
\core{令$G(u)=\int_0^u g(s)ds$。对$x\ne0$，
\[\int_0^{x^2}g(xt)dt=\frac{G(x^3)}x.\]
因此$f'(x)=3xg(x^3)-G(x^3)/x^2$。在0处必须独立用定义；连续性给$G(u)=g(0)u+o(u)$。}
\process{\begin{enumerate}\item 令$u=xt$，同时换微元和上下限，保留积分方向。\item 对$x\ne0$的商式求导，只使用连续性可保证的基本定理。\item 用$[f(x)-f(0)]/x$求分界点导数。\item 再比较$\lim f'(x)$与$f'(0)$。\end{enumerate}}
\fast{直接写$\int t g'(xt)dt$会擅自提高题设光滑性。换元是这类题的关键决策；$x<0$时同样按有向积分成立。}
\budget{中；12分题9—12分钟。}
\fullquestion{math2-2026-18}
\apply{$f(0)=0$；由连续性，$f(x)=g(0)x^2+o(x^2)$，故$f'(0)=0$。对非零$x$用上框公式，$3xg(x^3)\to0$，$G(x^3)/x^2=x[g(0)+o(1)]\to0$，所以$f'$在0连续。全过程不要求$g$可导。}
\related{2020-16：$\int_0^1f(xt)dt=G(x)/x$；2018-16：先将卷积积分改写成$\int_0^x(x-u)f(u)du$再求导。}

\section{原函数拼接}
\recognize{分段函数求“一个原函数”，选项各段只差常数。分别求导检验之外，还要在分界点把函数接起来。}
\core{原函数$F$须在整个区间满足$F'=f$，所以$F$一定连续。若$f$在分界点连续，各段原函数只需选择常数使$F$连续，即可在该点可导且导数为$f$。若$f$有跳跃间断，则不能在跨越该点的区间上有原函数（导数具有介值性）。}
\process{\begin{enumerate}\item 每段先验$F'=f$，快速排除积分公式或符号错误。\item 比较分界点左右极限，调整积分常数差。\item 检查该点差商或使用$f$连续的结论。\end{enumerate}}
\fast{各段常数并非互不相关；“加同一个$C$”表示整体常数自由，段与段的常数差则由连续拼接确定。}
\budget{易至中；选择题1—2分钟。}
\fullquestion{math2-2023-02}
\apply{左段原函数$\ln(x+\sqrt{1+x^2})+C_1$，右段$(x+1)\sin x+\cos x+C_2$。在0处左极限$C_1$、右极限$1+C_2$，故$C_1=1+C_2$，对应D。左段根式形式先辨识导数符号，再做常数匹配。}
\related{2016-02：$2(x-1)$与$\ln x$拼接；2014-10：周期奇函数先降到基本区间，再用原函数求函数值。}

\section{有理函数积分}
\recognize{分母已分解为一次因子和不可约二次因子，或分子接近分母导数。按因子拆分，再将二次项的分子配成导数与常数。}
\core{一次因子$(x-a)^m$须列$A_1/(x-a)+\cdots+A_m/(x-a)^m$；二次因子配一次分子。
\[\int\frac{q'(x)}{q(x)}dx=\ln|q(x)|+C,\quad
\int\frac{dx}{(x-a)^2+b^2}=\frac1b\arctan\frac{x-a}b+C\quad(b>0).\]}
\process{\begin{enumerate}\item 假分式先长除；真分式按实因子列部分分式。\item 清分母，代根求部分系数，再比较系数。\item 二次分母配方，分子配成分母导数。\item 定积分最后统一代上下限，核对对数与反正切系数。\end{enumerate}}
\fast{出现$(x-1)^2+1$直接想到反正切；出现$2x-2$想到对数。不要把不可约二次因子的分子只设成常数。}
\budget{中；填空2—4分钟，10分题7—10分钟。}
\fullquestion{math2-2025-17}
\apply{\[\frac1{(x+1)(x^2-2x+2)}=\frac1{5(x+1)}+\frac{-x+3}{5(x^2-2x+2)}.\]
把$-x+3=-\tfrac12(2x-2)+2$，积分立即分成两个对数和一个反正切。代0、1得$3\ln2/10+\pi/10$。选择这一路线由分母“线性乘二次”的结构决定。}
\related{2019-16：重一次因子须列完整幂次；2022-13：单二次分母直接配导数；2018-11：先部分分式再取无穷上限。}

\section{根式三角换元}
\recognize{分母有$\sqrt{x(1-x)}$、$\sqrt{a^2-x^2}$，或根式与反三角函数配套。优先让根式与微元相消。}
\core{\[x=a\sin t\Rightarrow\sqrt{a^2-x^2}=a\cos t,\quad
x=\sin^2t\Rightarrow\sqrt{x(1-x)}=\sin t\cos t\ (0\le t\le\pi/2).\]
根式$\sqrt{x^2+a^2}$可用$x=a\tan t$；$\sqrt{x^2-a^2}$可用$x=a\sec t$，并按区间处理绝对值。}
\process{\begin{enumerate}\item 同时观察根式、反三角函数和积分区间。\item 选择使反三角变成$t$、根式抵消$dx$的换元。\item 连同上下限一起换，不混用$x,t$。\item 对简化后的多项式或三角式积分。\end{enumerate}}
\fast{只消根式却使另一部分更复杂的换元未必合算。若题中有$\arcsin\sqrt x$，$x=\sin^2t$比$x=\sin t$更直接。}
\budget{易至中；选择题1—3分钟。}
\fullquestion{math2-2020-03}
\apply{完整题干中的$\arcsin\sqrt x/\sqrt{x(1-x)}$提示$x=\sin^2t$。$dx=2\sin t\cos t\,dt$与分母完全相消，原积分变为$2\int_0^{\pi/2}t\,dt=\pi^2/4$，选A。若先分部，会产生更复杂的根式导数。}
\related{2024-18第二问：ODE解代入后，再处理圆根式积分；2014-09：配方后的反正切标准型。}

\section{分部消积分}
\recognize{乘积中有对数、反三角函数、变限积分，另一因子易积分。优先让求导更简单的部分作为$u$。}
\core{\[\int u\,dv=uv-\int v\,du,\quad
\int_a^b p(x)F(x)dx=[P(x)F(x)]_a^b-\int_a^b P(x)F'(x)dx.\]
若$F$是无初等原函数的积分，分部可将$F$变回被积函数。选$P$时可加常数，使某端点边界项消失。}
\process{\begin{enumerate}\item 选$u=\ln(\cdot)$、反三角或变限积分；$dv$选容易积分的部分。\item 先计算边界项能否为0，再写剩余积分。\item 剩余式出现复合导数时立即换元。\item 反常积分须单独验证边界项极限。\end{enumerate}}
\fast{变限积分乘$x$，优先对$x$积分，避免先求变限积分本身。对于$e^{ax}\sin bx$，连续分部两次后解出原积分。}
\budget{中；填空2—4分钟，大题7—10分钟。}
\fullquestion{math2-2019-13}
\apply{令$F(x)=\int_1^x\sin(t^2)/t\,dt$，则$f=xF$。取$dv=x\,dx$，边界$[x^2F(x)/2]_0^1=0$；余下$-\frac12\int_0^1x\sin(x^2)dx$，换元$x^2$后得$(\cos1-1)/4$。读到外层$x$恰能补足内层$1/x$，即可预见分部后的简化。}
\related{2017-11：$\ln(1+x)/(1+x)^2$，边界项趋0；2018-15：反正切求导后再根式换元。}

\section{对称与平移}
\recognize{积分区间关于0对称，或题给$f(x+T)-f(x)$。先利用对称、换元和平移，通常可减少一半计算或避开未知函数表达式。}
\core{\[\int_{-a}^a f=\int_0^a[f(x)+f(-x)]dx,\quad
\int_a^b f(x)dx=\int_a^b f(a+b-x)dx.\]
若$f(x+T)-f(x)=h(x)$，则长度相同的区间之差可用已知$h$积分；跨越端点的部分先拆开。}
\process{\begin{enumerate}\item 看区间中心、被积函数奇偶性与平移量。\item 通过$x\mapsto-x$或$x\mapsto x+T$写第二个积分。\item 将未知函数部分与已知总积分组合消去。\item 最后只计算简单的修正项。\end{enumerate}}
\fast{有$e^x$不意味着没有对称性；可先看$f(x)+f(-x)$。未知$f$加上平移关系时，不必先求$f$，因为可能存在任意周期部分。}
\budget{易至中；填空2—3分钟。}
\fullquestion{math2-2023-15}
\apply{把$[1,3]$拆成$[1,2]$与$[2,3]$。后者令$x=u+2$，由$f(u+2)=f(u)+u$得到
\[\int_1^3f=\int_1^2f+\int_0^1[f(u)+u]du=\int_0^2f+\frac12=\frac12.\]
触发点是新旧区间仅有一段平移2，而不是要求解函数方程。}
\related{2018-05：对称区间先取偶部；2012-04：分周期平移，把不同波瓣配对比较。}

\section{积分大小比较}
\recognize{问多个积分的排序、正负，函数看似无法求原函数。区间相同时比较被积函数，区间不同时先拆分、平移或统一参数。}
\core{若$f\le g$且均可积，则$\int f\le\int g$；连续且在某个内点严格小于，通常可推出积分严格小于。常用
\[\frac{x}{1+x}\le\ln(1+x)\le x\quad(x\ge0),\qquad\sin x<x\quad(x>0).\]
与固定分母比较时，还要记录分母的正负。}
\process{\begin{enumerate}\item 将比较对象放到同一积分区间。\item 先找共同分母或共同正因子，只比较剩余部分。\item 基本不等式不足时，对两函数之差求导。\item 用桥梁函数串起大小链，无需精确积分值。\end{enumerate}}
\fast{桥梁函数$x$、$x/2$、常数1常比直接两两作差更省。积分为0的奇部先去掉，再比较剩余偶部。}
\budget{中；选择题2—3分钟。}
\fullquestion{math2-2022-07}
\apply{在$0<x<1$，$\ln(1+x)>x/2$，因为其与$x/2$之差导数$1/(1+x)-1/2>0$，所以$I_2>I_1$。又$\ln(1+x)<x$且$1+\cos x>1$，故$I_2<\int_0^1x\,dx$；$2/(1+\sin x)>1$给$I_3>\int_0^1x\,dx$，所以$I_1<I_2<I_3$，选A。}
\related{2012-04：指数权重使后续负波瓣更重；2018-05：先对称再比较。}

\section{反常积分定界}
\recognize{积分到无穷，或端点分母为0，要求参数范围。先列全部瑕点，对每个端点分别求局部等价式，再取条件交集。}
\core{\[\int_0^1x^{-p}dx\text{收敛}\iff p<1,\quad
\int_1^\infty x^{-p}dx\text{收敛}\iff p>1;\]
\[\int_A^\infty\frac{dx}{x(\ln x)^q}\text{收敛}\iff q>1\quad(A>1).\]
正函数比值趋正有限常数则同敛散；符号恒负时可对绝对值判断。无穷区间与瑕点方向的阈值相反。}
\process{\begin{enumerate}\item 列出端点和内部异常点，在普通点拆开。\item 每一端只保留主导因子，确定局部幂次。\item 非临界幂次直接比较；幂次恰为1再看对数。\item 所有端点条件求交集，临界参数单独验。\end{enumerate}}
\fast{在$x\to1$，$\ln x\sim x-1$，能与分母的$(1-x)$抵消；不能把它当成常数。在$x\to0$，$\ln x$虽无界却慢于任何负幂。}
\budget{中；选择填空2—4分钟。}
\fullquestion{math2-2022-05}
\apply{原积分被积函数为$\ln x/[x^p(1-x)^{1-p}]$。在0附近以$|\ln x|/x^p$判断，收敛须$p<1$；在1附近，$\ln x\sim-(1-x)$，故等价于$-(1-x)^p$，收敛须$p>-1$。取交集$(-1,1)$，选A。}
\related{2013-04：有限瑕点与无穷端同时控制；2026-11：0附近$\arctan x\sim x$，无穷处$\arctan x\to\pi/2$；2023-06：先令$u=\ln x$确定定义域，再讨论参数极值。}

\section{审敛与反例}
\recognize{“平方可积是否推出可积”“收敛是否存在幂次衰减”是性质命题。先用$p$型和对数临界型试方向，再用比较判别法证明真命题。}
\core{无穷区间上，$\int f^2<\infty$不能推出$\int f<\infty$；有限区间上可由柯西不等式推出绝对可积。若非负$f$满足$x^pf(x)$有有限极限且$p>1$，则尾部$f(x)\le C/x^p$，积分收敛。反之收敛不必具有任何$p>1$的幂次衰减。}
\process{\begin{enumerate}\item 分清积分区间是否无穷，命题是充分还是必要。\item 用$f=(1+x)^{-a}$检验平方与一次的差异。\item 幂次结论过强时，试$1/[x(\ln x)^2]$的平移版。\item 真命题写出尾部上界并给比较积分。\end{enumerate}}
\fast{反常积分穿过内部瑕点时左右必须各自收敛；$\int_{-1}^1dx/x$不能靠奇函数抵消宣布收敛。对称主值与通常反常积分不同。}
\budget{中；多命题选择题3—4分钟。}
\fullquestion{math2-2024-07}
\apply{①取$f=1/(1+x)$，平方积分收敛而一次发散。②有限极限给$x^pf(x)$尾部有界，与$1/x^p$比较，命题真。③取$f=1/[(x+2)\ln^2(x+2)]$，积分收敛，但任意$p>1$均有$x^pf(x)\to\infty$，命题假。仅一个正确，选B。}
\related{2019-03：内部瑕点不能左右抵消；2021-11：偶函数先化为两倍半轴积分再换元；2025-11：先计算参数反常积分再求参。}
